\documentclass[12pt]{article}
\usepackage[margin=1in]{geometry}
% \usepackage[UTF8,fontset=fandol]{ctex}
\usepackage{amsmath,amssymb,graphicx,booktabs,tabularx,array,setspace,float}
\usepackage[colorlinks=true,linkcolor=blue,citecolor=blue,urlcolor=blue]{hyperref}
\setstretch{1.1}
\setlength{\parindent}{2em}
\renewcommand{\figurename}{Figure}
\renewcommand{\tablename}{Table}

\title{Directional Congestion and Tidal Lanes}
\author{
Yige Duan\thanks{Antai College of Economics and Management, Shanghai Jiao Tong University; \href{mailto:yige.duan@sjtu.edu.cn}{yige.duan@sjtu.edu.cn}.}
\and
Xiaoruo Feng\thanks{Peking University HSBC Business School; \href{mailto:xiaoruo.feng@stu.pku.edu.cn}{xiaoruo.feng@stu.pku.edu.cn}.}
\and
Yizhen Gu\thanks{Peking University HSBC Business School; \href{mailto:yizhengu@phbs.pku.edu.cn}{yizhengu@phbs.pku.edu.cn}.}
}
\date{March 29, 2026}

\begin{document}
\maketitle

\begin{abstract}
This paper studies directional congestion and the potential effects of tidal-lane policies in large cities. The analysis combines high-frequency road-speed data and fine-grid commuting flows for Beijing. The paper asks whether directional asymmetry is an important feature of urban congestion and how policies that reallocate road capacity across directions should be evaluated. The preliminary evidence shows substantial peak-hour directional imbalance in travel speeds together with a pronounced spatial mismatch between residences and jobs. To interpret these patterns, the paper develops a quantitative spatial-equilibrium model with directional commuting costs, route assignment, and endogenous congestion.
\end{abstract}

\section{Introduction}
Urban transportation infrastructure is costly to build and slow to expand, especially in dense central-city environments where land is scarce and construction is disruptive. Yet congestion on existing road networks is severe. In large Chinese cities, peak-hour traffic substantially reduces effective travel speeds, imposes large time losses on commuters, and lowers the efficiency with which existing infrastructure is used. Because these losses affect both accessibility and the allocation of economic activity across space, congestion is not only a transportation problem; it is also a problem of economic efficiency and urban welfare.

Expanding road infrastructure is costly and slow, governments need policies that make more intensive use of the existing network. This paper focuses on one particularly relevant manifestation of urban congestion: directional scarcity. In large monocentric cities, morning inflows toward employment centers and evening outflows toward residential areas place unequal pressure on the two directions of the same physical corridor. The resulting asymmetry is directly relevant for policy because the scarcity at stake is not only how much road capacity the city has in total, but also how that capacity is allocated across directions and times of day.

We study this issue through the lens of \emph{tidal lanes}, that is, policies that reallocate existing lane capacity toward the peak-flow direction. Relative to large-scale road expansion, tidal lanes are attractive because they can be implemented more quickly, at lower fiscal cost, and in a more targeted way on corridors with strong directional imbalance. But evaluating such policies requires more than measuring local changes in speed. A directional capacity change can alter route substitution, bilateral commuting costs, home-workplace matching, and the spatial allocation of residents and jobs. In other words, the relevant policy question is not only whether tidal lanes reduce congestion on treated segments in the short run, but also whether they improve system-wide efficiency and welfare once network spillovers and spatial adjustment are taken into account.

Against this background, the paper asks three related questions. First, how important is directional asymmetry in urban road congestion once we measure travel conditions on a directed network? Second, which corridors appear to be plausible candidates for capacity reallocation policies such as tidal lanes? Third, how should such policies be evaluated when commuting costs are directional and congestion feeds back into route choice and spatial allocation?

The paper makes two related contributions. On the measurement side, it constructs a directed road network that recovers direction-specific travel conditions on the same centerline corridor and combines those conditions with grid-level commuting flows. On the modeling side, it embeds the resulting ordered commuting costs in a quantitative spatial framework with directional commuting frictions and endogenous congestion. Together, these pieces provide a framework for studying whether road-capacity reallocation improves urban traffic performance and welfare.

\section{Data and Variables}
\subsection{Data Sources and Constructed Objects}
The empirical analysis uses two high-resolution data sources for Beijing. The first is a road-speed dataset with segment-level observations at high temporal frequency, including geometry, road class, and direction information. The second is a fine-grid origin-destination commuting dataset that records home-workplace flows across small spatial cells.

These data are combined in three steps. First, raw road geometry is regularized into an undirected centerline network and then duplicated into two directed links, AB and BA. Second, raw speed records are matched to directed centerlines and aggregated into direction-specific speed panels by hour and peak period. Third, commuting flows are aggregated to a square-grid system and mapped to the directed road graph so that each origin-destination pair can be assigned a direction-sensitive travel cost.

The resulting empirical objects are: an undirected centerline network, a directed AB/BA network, a centerline-level directional speed panel, grid-level resident and employment distributions, and an ordered commuting-cost matrix defined over the grid system.

\subsection{Spatial Distribution of Residents and Jobs}
The spatial distribution of residents and jobs is the first descriptive motivation for directional congestion. Employment is more concentrated than residence, implying systematic morning inflows toward central employment clusters and reverse flows later in the day. In the figure, each grid is shaded by the percentile rank of residents or jobs so that the map emphasizes relative spatial concentration rather than absolute population levels.

\begin{figure}[H]
    \centering
    \includegraphics[width=0.92\textwidth]{todo_outputs/figs/figure1_spatial_distribution_residents_jobs.pdf}
    \caption{Spatial distribution of residents and jobs in the Beijing square-grid system. Each grid is colored by the percentile rank of residents or jobs, with no grid borders.}
    \label{fig:spatial-rj}
\end{figure}

\subsection{Road Speeds and Directional Asymmetry}
The directed speed panel contains more than 24 million matched observations after preprocessing and aggregation. The central empirical object is the difference in peak-period travel conditions across directions on the same physical corridor.

Directional asymmetry is measured by comparing AB and BA speeds on the same centerline within a given time interval. Let $\bar v_{c,d,\tau}$ denote the mean observed speed on centerline $c$ in direction $d \in \{AB,BA\}$ during interval $\tau$, where the mean is taken over all speed records assigned to that centerline, direction, and interval. We define the asymmetry ratio as
\[
a_{c,\tau} = 1 - \frac{\min(\bar v_{c,AB,\tau},\bar v_{c,BA,\tau})}{\max(\bar v_{c,AB,\tau},\bar v_{c,BA,\tau})}.
\]
This measure is zero when the two directions have the same average speed over the interval and rises as directional imbalance becomes stronger. Figure~\ref{fig:asym-dist} shows the distribution of this statistic in the AM and PM peaks. The distributions are shifted well above zero, especially in the morning peak, indicating that directional congestion is a first-order feature rather than incidental noise.

\begin{figure}[H]
    \centering
    \includegraphics[width=0.82\textwidth]{../data_work/outputs/match_projection_complete_v7_latest/figures/hist_speed_ratio_am_pm.png}
    \caption{Distribution of directional asymmetry in the AM and PM peaks.}
    \label{fig:asym-dist}
\end{figure}

The asymmetries documented above are not merely local curiosities. In the sample, directional imbalance is sufficiently strong and persistent to make capacity reallocation a plausible policy margin. This is precisely the environment in which tidal-lane interventions may matter: when demand is directionally unbalanced and bilateral commuting costs differ meaningfully by direction.

\subsection{Summary Statistics}
Table~\ref{tab:summary} reports the main descriptive statistics used in the draft. At the grid level, the average cell contains 2,372 residents and 2,372 workers, but the medians are much smaller, reflecting substantial spatial concentration. At the segment level, the average centerline length is about 1.47 km, the mean estimated lane count is slightly above 3, and the average asymmetry ratio is about 0.195.

\begin{table}[!htbp]
\centering
\caption{Summary Statistics}
\label{tab:summary}
\begin{tabular}{llrrr}
\toprule
Panel & Variable & Mean & SD & Median \\
\midrule
Panel A. Grid-level & Residents & 2,372.079 & 6,556.008 & 74.000 \\
 & Workers & 2,372.079 & 8,468.076 & 61.000 \\
 & Distance to city center (km) & 77.539 & 35.958 & 76.171 \\
Panel B. Segment-level & Length (m) & 389.782 & 1,585.561 & 7.071 \\
 & Number of lanes & 3.175 & 1.177 & 3.615 \\
 & Average speed (km/h) & 37.335 & 11.874 & 34.397 \\
 & Morning peak speed (km/h) & 35.396 & 12.342 & 32.371 \\
 & Evening peak speed (km/h) & 33.958 & 12.596 & 31.145 \\
 & Asymmetry ratio & 0.192 & 0.174 & 0.131 \\
\bottomrule
\end{tabular}
\end{table}


\section{Model}
\subsection{Directed Commuting Costs}
The model partitions the city into grid cells indexed by $i,j = 1,\dots,N$. Let $L_i$ denote residents in grid $i$, $H_j$ jobs in grid $j$, and $M_{ij}$ the commuting flow from residence $i$ to workplace $j$. Let $t_{kl}$ denote travel time on directed edge $(k,l)$ of the road graph, and let $\tau_{ij}$ denote the bilateral commuting cost from $i$ to $j$.

The key object is an ordered commuting-cost matrix:
\[
\tau_{ij} = \min_{p \in \mathcal{P}_{ij}} \sum_{(k,l)\in p} t_{kl},
\]
where $\mathcal{P}_{ij}$ is the set of feasible directed paths from $i$ to $j$. Because travel times are directional, in general $\tau_{ij} \neq \tau_{ji}$. This is the central departure from standard symmetric commuting-cost formulations.

\subsection{Commuting Flows and Spatial Allocation}
The model follows a quantitative spatial structure with residential amenity and workplace productivity fundamentals. Define
\[
U_i^\theta = \bar u_i^\theta (L_i/\bar L)^{\beta\theta},
\qquad
A_j^\theta = \bar a_j^\theta (H_j/\bar L)^{\alpha\theta},
\]
where $\bar L = \sum_i L_i$ is total population, $\bar u_i^\theta$ is a residential fundamental, and $\bar a_j^\theta$ is an employment fundamental. Commuting mass between $i$ and $j$ is
\[
m_{ij} = \tau_{ij}^{-\theta} U_i^\theta A_j^\theta,
\]
and the associated commuting flow is
\[
M_{ij} = \bar L \cdot \frac{m_{ij}}{\sum_{r,s} m_{rs}}.
\]
Equilibrium residents and jobs satisfy
\[
L_i = \sum_j M_{ij},
\qquad
H_j = \sum_i M_{ij}.
\]

This structure makes the economic role of directional commuting costs transparent. Any change in the ordered matrix $\tau_{ij}$ shifts bilateral commuting masses, which in turn changes residence-workplace matching and the implied spatial distribution of population and employment. In that sense, the model links directional road performance to both short-run traffic outcomes and longer-run spatial allocation.

\subsection{Baseline Inversion of Fundamentals}
The model takes observed residents $L_i^{obs}$, jobs $H_j^{obs}$, and measured commuting costs $\tau_{ij}$ as inputs and inverts the residential and employment fundamentals so that the baseline equilibrium is consistent with the observed spatial distribution. Let
\[
\ell_i^{L,obs} = \frac{L_i^{obs}}{\bar L},
\qquad
\ell_j^{H,obs} = \frac{H_j^{obs}}{\bar L}.
\]
Then the inverted fundamentals satisfy the system
\[
\bar u_i^\theta
=
\frac{(\ell_i^{L,obs})^{1-\beta\theta}}
{\sum_j \tau_{ij}^{-\theta}\bar a_j^\theta(\ell_j^{H,obs})^{\alpha\theta}},
\]
\[
\bar a_j^\theta
=
\frac{(\ell_j^{H,obs})^{1-\alpha\theta}}
{\sum_i \tau_{ij}^{-\theta}\bar u_i^\theta(\ell_i^{L,obs})^{\beta\theta}}.
\]
This inversion step is important because it allows the baseline model to start from an empirically disciplined city rather than from arbitrary fundamentals.

\subsection{Route Assignment and Congestion}
Given commuting flows, route assignment is all-or-nothing along shortest paths. If $p_{ij}^*$ denotes one shortest path between $i$ and $j$, then edge flow on directed edge $(k,l)$ is
\[
\phi_{kl} = \sum_{i,j} M_{ij}\mathbf{1}\{(k,l)\in p_{ij}^*\}.
\]

Let $n_{kl}$ denote lane capacity on edge $(k,l)$ and let $t_{kl}^{ff}$ denote free-flow travel time. Congested travel time is given by
\[
t_{kl} = t_{kl}^{ff} \left(\frac{\phi_{kl}}{n_{kl}}\right)^\lambda,
\]
where $\lambda$ is the congestion elasticity. This is the margin through which tidal-lane policies enter the model: reallocating lane capacity changes $n_{kl}$ for treated directions, which changes edge travel times, shortest-path commuting costs, and finally the equilibrium allocation of commuting flows, residents, and jobs.

\subsection{Model Results}
\textbf{Baseline fit.} TBD.

\textbf{Directional versus symmetric commuting costs.} TBD.

\textbf{Congestion and route-assignment implications.} TBD.

\subsection{Counterfactual Design}
The quantitative framework is used for several policy and benchmarking exercises.

\textbf{Counterfactual 1: symmetric versus directional commuting costs.} TBD.

\textbf{Counterfactual 2: adding tidal lanes to the most asymmetric corridors.} TBD.

\textbf{Counterfactual 3: adding road capacity and the implications for the fundamental law of congestion.} TBD.

\textbf{Counterfactual 4: welfare and distributional comparisons across scenarios.} TBD.

\end{document}
